\section{与 Project 0 精确解程序的连接}

为了检查数值解是否合理，Project 2 在输出数值解之后，会调用之前 Project 0 中写好的
Riemann 精确解程序。这样做的好处是不用在 Project 2 里再写一遍精确解算法，
两个程序各自负责一件事。

\subsection{调用方式}

Project 0 的可执行文件是
\[
\texttt{\_Analysical\_Solution\_Solver\textbackslash riemann\_exact.exe}.
\]
Project 2 把左右状态、\(\gamma\)、计算域、网格点数、输出时间和输出文件名都拼成命令行参数，
然后通过
\begin{lstlisting}
system(command);
\end{lstlisting}
启动精确解程序。

如果精确解程序不存在，Project 2 会给出编译提示；如果外部程序返回错误，
Project 2 也会停止当前计算。

\subsection{为什么要传同一个网格}

Project 2 会把自己的
\[
x_{\min},\quad x_{\max},\quad x_0,\quad N_x,\quad \gamma
\]
传给 Project 0。这样 Project 0 生成的 exact 文件和 numerical 文件在同一组
\[
x_i=x_{\min}+i\frac{x_{\max}-x_{\min}}{N_x-1}
\]
上取值。后面比较误差时就可以直接逐点比较，不需要再做插值。

\subsection{路径和引号}

因为是在 Windows 下运行，路径中可能有空格，所以命令字符串里需要对可执行文件和输出文件名加引号。
这一步主要是为了保证 \texttt{system()} 能正确找到外部程序。
